\chapter{Model editing과 unlearning: 좁게 바꾸고 정확히 되돌리기}
\label{bg:chapter-editing-unlearning}

Fine-tuning은 dataset 분포 전체에 맞춰 model behavior를 바꾸지만, 때로는 ``A의 CEO는 B가 아니라
C다''처럼 몇 개 fact만 빨리 수정해야 한다. 반대로 삭제 요청은 특정 data의 영향을 제거해야 한다.
두 문제 모두 update 범위, 부작용, provenance, rollback을 요구하지만 성공 기준은 다르다.

\section{Model editing의 네 평가축}

\concept{model-editing}{모델 편집}은 특정 edit request $e=(x_e,y_e)$를 반영하면서 나머지 behavior를
보존하려는 수정이다. ROME은 transformer MLP의 factual association을 찾아 rank-one update를 적용하고
\parencite{meng2022rome}, MEMIT은 이를 많은 memory edit로 확장했다
\parencite{meng2022memit}. 이들은 ``fact가 한 주소에 완전히 저장된다''는 증명이 아니라, 특정
benchmark와 architecture에서 효과적인 intervention을 보인 방법이다.

평가는 최소 네 축으로 나눈다.

\begin{enumerate}
  \item \textbf{efficacy}: edit prompt에서 새 답을 내는가?
  \item \textbf{paraphrase/generalization}: 표현이 달라도 같은 correction을 쓰는가?
  \item \textbf{specificity}: 비슷하지만 다른 entity·relation은 보존되는가?
  \item \textbf{retention}: unrelated capability와 이전 edit가 유지되는가?
\end{enumerate}

\concept{locality}{국소성}은 세 번째와 네 번째를 포괄한다. 단순하게 edit prompt log-probability만
높이면 model이 주변 사실까지 바꾸거나 같은 token을 과도하게 선호할 수 있다. Sleep batch에는 edit
example뿐 아니라 neighborhood, counterfactual, old capability probe가 필요하다.

\section{많은 edit는 다시 continual learning과 capacity 문제가 된다}

한두 edit는 저랭크 update로 가능해 보여도 수천·수백만 edit를 누적하면 interference가 늘어난다.
같은 subject의 시간별 fact가 충돌하고, 이전 edit가 새 edit에 의해 덮이며, 어느 update를 제거해야 하는지
모호해진다. 매일 editing을 반복하는 시스템은 결국 다음 중 하나를 선택해야 한다.

\begin{itemize}
  \item latest fact만 weights에 남기고 history는 temporal external memory에 둔다.
  \item edit를 versioned adapter에 격리하고 주기적으로 merge/distill한다.
  \item 반복 조회되고 안정적인 fact만 parametric promotion하고 나머지는 retrieval한다.
  \item global knowledge refresh는 큰 offline model release에 맡긴다.
\end{itemize}

즉 editing은 sleep-time compute의 가능한 operator이지만, 무한 append 가능한 장기기억 자체는 아니다.

\section{Unlearning은 delete key보다 강한 요구다}

\concept{unlearning}{머신 언러닝}은 지정 data가 학습에 없었던 counterfactual model에 가까워지도록
그 영향을 제거하는 절차다. Sharded, Isolated, Sliced, Aggregated(SISA) training은 data shard별 model을
분리해 삭제 시 영향받은 부분만 재학습하는 방향을 제시했다 \parencite{bourtoule2021unlearning}.

External memory는 item pointer를 지우기 쉬워도 cache, embedding replica, graph edge, summary,
downstream training dataset까지 삭제가 전파되어야 한다. Parametric memory는 어느 weight가 한 record에
기여했는지 역함수가 없으므로 exact deletion이 훨씬 어렵다. ``답을 더 이상 말하지 않는다''와 ``data
영향이 제거되었다''도 다르다.

\begin{table}[htbp]
\centering
\caption{수정·삭제 destination별 복구 성질}
\label{tab:bg-edit-unlearn}
\small
\begin{tabularx}{\textwidth}{p{0.18\textwidth} X X X}
\toprule
Destination & 수정 단위 & 삭제/롤백 & 잔여 위험 \\
\midrule
raw text/log & record/segment & pointer+replica 삭제 & 파생 summary·dataset \\
vector index & vector+metadata & tombstone+rebuild & cache와 stale replica \\
knowledge graph & node/edge/time interval & edge version rollback & derived edge와 conflict \\
adapter/LoRA & file 또는 rank component & version detach & merged artifact, shared behavior \\
base weights & diffuse parameters & retrain/edit/unlearn & 비국소 영향, 검증 비용 \\
\bottomrule
\end{tabularx}
\end{table}

\section{Provenance graph가 architecture choice를 바꾼다}

권리 만료가 잦거나 user delete SLA가 짧다면 external memory나 isolated adapter의 운영 가치가 커진다.
Stable public knowledge처럼 많은 query가 반복되고 source set이 고정될 때는 global parametric refresh가
유리할 수 있다. 따라서 parametric/external 선택은 accuracy 비교만이 아니라 expected delete frequency,
scope, edit conflict, rollback time을 포함한 lifecycle cost 비교다.

\begin{cautionbox}[Negative answer만으로 unlearning을 증명하지 말 것]
모델이 특정 prompt에 답하지 않도록 fine-tune한 것은 suppression일 수 있다. membership inference,
paraphrase, related representation, downstream model lineage를 포함해 요구 수준을 정의해야 한다. 완전한
counterfactual retraining과 근사 unlearning을 보고서에서 구분한다.
\end{cautionbox}
